\section{Limitations and Open Questions}\label{sec:limits}

\paragraph{Access, not information flow.} A grant decides which actions an
agent may take. It does not track what the agent learns from an allowed action
or where that data goes next through other allowed actions: an agent allowed
to read a file and to reach a host can send the file's contents to that host.
Access control does not by itself enforce information-flow
policies~\cite{sabelfeld}; systems that track data flow through an agent's
tool calls~\cite{camel,fides} address a different question.

\paragraph{What the proofs cover.} The proofs are \translationtrust. The step
bound is proven for a counted model and is loose (\S\ref{sec:proofs}); the
measured cost follows the model's step count on the measured points, not
provably on every input\claim{deployed-cost-general}. The run-level budget
bound is proven for the grant signer's calls as \file{Signer.lean} models
them; that the contract's code is that model is
trusted\claim{wiring-budget-calls}.

\paragraph{The wiring is trusted, and has known gaps.} That the argument
mapping keeps distinct arguments distinct\claim{xdr-injective} and that the
relay's supplied and recorded trees cover every authorized
context\claim{trees-cover} are tested or argued, not proven. The host's path
resolution, URL parsing, redirect handling and payload decoding, and the
operating system calls beneath them, are tested only. Table~\ref{tab:gaps}
lists the gaps the paper states: under an empty grant the grant signer's own
address can authorize any call on another contract; a holder's signature
does not cover the grant id, so a key should hold one grant; no caveat binds
a signed transaction's fee, sequence number, time bounds or memo; network
caveats bind names, not the addresses they resolve to, and not the HTTP
method or path; and a file renamed after its descriptor was checked, or
hard-linked into the grant, is reached. The account's document must meet two
requirements nothing on chain checks\claim{document-requirements}; the grant
signer cannot check that the document names it, since it cannot see which
account delegated to it.

\paragraph{Scope of the hidden-call guarantee.} The grant signer sees the
contexts that need the account's authorization. Effects that need none, such
as a \code{transfer\_from} against a standing allowance or funds a contract
already holds for the account, are invisible to every caveat and budget, and
the paper assumes no allowed contract has such a
position\claim{no-standing-allowance}. The relay likewise cannot see calls
that ask for no authorization or calls a contract makes inside its own
\code{\_\_check\_auth}, and bounds the fee by a cap rather than a caveat.

\paragraph{Expressiveness.} Caveats are a closed set: membership of a field
in a list, an upper bound on one \code{i128} argument, and equality of one
argument. Each applies to every context of a tree alike, so a grant over two
contracts and two functions allows their cross product, and an argument
caveat reads the same position of every context (\S\ref{sec:grants}). There
is no caveat scoped to one contract and function, no lower bound, no expiry
and no caveat across two arguments. Allowances are cumulative and never reset:
the periodic and refillable allowances of session-key systems and Zodiac
Roles~\cite{zodiacRoles,smartSessions} have no counterpart, and would need a
time model in the budget code and its proofs. Budgets are enforced only at the
Soroban boundary, and each root grant meters one asset; transfers of other
assets are bounded only by caveats.

\paragraph{Keys and storage.} Holder keys at the Soroban boundary are ed25519
only. Grants live in persistent contract storage, extended to the network
maximum on every write; an archived grant fails authorization until it is
restored, which fails closed.

\paragraph{Evaluation.} The attack suite ran in Soroban simulation and at the
host, not on a local network and not through the relay; the relay has
functional tests only, with no cost, throughput or attack-suite measurement.
Host timings are the medians of one run on one machine.

\paragraph{Dfns.} Whether real Dfns's \code{Transaction} kind signs a
transaction whose source is a channel account\claim{dfns-transaction-source},
and whether its \code{Hash} kind signs 32 bytes directly as an ed25519
message\claim{dfns-hash}, are not stated by its reference and were not tried
against the service. The relay signs through \code{Hash}, which the testnet
run shows the network accepts when the signature is right.
